\section{Problem Formulation}

Let $\mathcal{S}=\{s_1,\ldots,s_N\}$ denote $N$ asynchronous video streams. At time $t$, each stream can provide cheap features $x_{i,t}$ such as motion activity, lightweight anomaly scores, object counts, or CLIP-style prompt similarity. The agent chooses a set of visual observation actions $A_t$ subject to budget $B_t$. An action may query a VLM on a stream segment, crop, or question:

\[
a_{i,t} = \texttt{look}(s_i, t, \Delta, r, q),
\]

where $\Delta$ is temporal extent, $r$ is resolution or region, and $q$ is an optional natural-language incident query. Each action has cost $c(a_{i,t})$ in GPU-seconds and VLM calls.

The action space is intentionally hierarchical. A system may first spend a cheap scan on every stream, then choose a small set of streams for medium-cost feature extraction, and finally call a VLM on a selected clip or crop. In this paper we evaluate the hardest bottleneck explicitly: the number of expensive VLM calls available at each timestep. The same protocol can also enforce a global episode-level GPU budget or memory budget.

The episode contains latent event intervals $\mathcal{E}=\{e_j\}$, where each event has a stream, start time, end time, label, severity, and language description. A successful monitoring agent must detect events quickly while controlling operational cost:

\[
\max_{\pi} \ \mathbb{E}\left[\sum_{e_j \in \mathcal{E}} u(e_j, \tau_j, \hat{d}_j) - \lambda_c \sum_t c(A_t) - \lambda_f F\right],
\]

where $\tau_j$ is detection delay, $\hat{d}_j$ is the generated incident description, and $F$ is the number of false alarms.

For evaluation, an event is counted as detected if the agent raises a positive alert during the event interval or within a fixed grace window after event onset. Duplicate alerts for the same event do not improve recall but do increase operational cost. A false alarm is any positive alert on a stream-time pair that does not overlap a ground-truth event. This event-centric accounting prevents a policy from scoring well by repeatedly inspecting the same easy event or by flooding the operator with alerts.

We report the following metrics:

\begin{itemize}
    \item \textbf{Event Recall@Budget}: fraction of ground-truth events detected under a fixed VLM-call or GPU-time budget.
    \item \textbf{Time-to-Detect}: delay between event start and first correct alert.
    \item \textbf{False Alarms per Hour}: number of positive alerts without an active event, normalized by monitored time.
    \item \textbf{GPU-Seconds per Hour}: deployment-oriented compute cost.
    \item \textbf{VLM Calls per Event}: expensive perception usage normalized by event count.
    \item \textbf{Incident Summary Correctness}: agreement between generated natural-language reports and event annotations or human judgments.
\end{itemize}

This formulation makes perception allocation measurable. A model that answers accurately after receiving the right clip is insufficient if it cannot decide which clip to inspect.
